\chapter{The evidence model}\label{ch:evidence-model}

Every later chapter makes quantitative statements about LATAM Bank's data: a share of rows, a rate of events, a
difference between the main and backup copies, a gain of a model over a benchmark. This chapter fixes what such a
statement means, which uncertainty it must carry, and what it may and may not be used for. It sets out the estimands
and the unit of analysis (\cref{sec:em-estimands}), the claims ledger that every number in the report comes from
(\cref{sec:em-ledger}), the interval estimators and the rule that chooses among them
(\cref{sec:em-iid,sec:em-clusters,sec:em-few,sec:em-series}), the control of false discoveries across many tests
(\cref{sec:em-multiplicity}), what a registered re-estimation can and cannot establish
(\cref{sec:em-registered}), and the rule that turns a model's out-of-time gain into a readiness verdict
(\cref{sec:em-gate}).

\section{Estimands and the unit of analysis}\label{sec:em-estimands}

The dataset records events: transactions, digital events, contacts, sends, complaints. Events are not independent
draws. A customer makes many transactions, an agent handles many contacts, a campaign sends many messages, and the
days of a calendar follow one another. We write the data as a nested structure.

\begin{definition}[Clustered observations]\label{def:clusters}
Let $g = 1, \dots, G$ index clusters (customers, agents, branches, campaigns or days) and $i = 1, \dots, m_g$ the
observations of cluster $g$, with $n = \sum_g m_g$. An observation carries an outcome $y_{gi}$ and, for a ratio, a
size $d_{gi} > 0$ (for a proportion, $d_{gi} = 1$ and $y_{gi} \in \{0, 1\}$).
\end{definition}

The estimand of most descriptive claims is a ratio of population totals,
\begin{equation}\label{eq:ratio}
  \theta = \frac{\E \bigl[\sum_i y_{gi}\bigr]}{\E \bigl[\sum_i d_{gi}\bigr]},
  \qquad
  \hat\theta = \frac{\sum_{g} Y_g}{\sum_{g} D_g},
  \quad Y_g = \sum_{i} y_{gi},\; D_g = \sum_{i} d_{gi},
\end{equation}
a share, a rate per exposure or a mean per unit. The estimator $\hat\theta$ is the same whatever the dependence
between observations; its sampling distribution is not. The independent unit is the cluster, not the row, and the
interval of a claim must be computed over the unit that is actually independent.

\begin{assumption}[Independence between clusters]\label{ass:clusters}
The cluster totals $(Y_g, D_g)$, $g = 1, \dots, G$, are independent and identically distributed. Observations within
a cluster may be arbitrarily dependent.
\end{assumption}

\Cref{ass:clusters} is weaker than independence of rows, which is the special case $m_g = 1$. It is the assumption
behind every clustered interval in this report; where it is doubtful, because clusters share a calendar day or a
branch, the chapter that uses it says so.

\section{The claims ledger}\label{sec:em-ledger}

No number derived from data is typed into this report. Every such number is a field of a row of the claims ledger,
and the text cites it with a macro that the build generates from that row.

\begin{definition}[Claim]\label{def:claim}
A claim is a row with an identifier, a kind, a status, the path and the \shaname{} digest of the artefact it was
computed from, and the fields its kind requires: a count needs a value; a proportion needs a numerator, a positive
denominator, an interval, the method that produced the interval and the cluster unit; a test needs its statistic, its
p-value, its family and the false-discovery adjusted value $q$ of \cref{sec:em-multiplicity}.
\end{definition}

The build fails when a chapter cites a claim that does not exist, cites a field the claim does not have, or contains a
typed number in prose or in mathematics. A claim's status records what it may be used for.

\begin{description}[style=nextline]
  \item[Exploratory] A number from the notebooks or from the result tables they wrote. The notebooks searched among
  many specifications, and some of their intervals treat rows as independent; an exploratory claim describes what
  was found, and an interval computed under row independence is marked with \naivemark{} wherever it is printed.
  \item[Registered] A re-estimation under the specification frozen before it ran (\cref{sec:em-registered}). It is
  the only status that may decide whether a readiness gate's verdict holds.
  \item[Positive control] A number measured on data with a planted, known answer, such as anomalies injected to
  measure a detector's recall. It shows whether a method can find what it should; it is never evidence about the bank
  and never enters a table of real-data results.
\end{description}

Whole tables are not split into claims. A table is rendered from one stored result, and its caption prints the path
and the first characters of the \shaname{} digest of that result, so a reader can check that the table and the file
agree.

\section{Intervals under independence}\label{sec:em-iid}

When rows are independent, the share $\hat p = Y/n$ of $Y$ successes in $n$ trials has the Wilson interval
\autocite{wilson1927}, obtained by inverting the score test of $H_0\colon p = p_0$:
\begin{equation}\label{eq:wilson}
  \frac{\hat p + \frac{z^2}{2n}}{1 + \frac{z^2}{n}}
  \;\pm\;
  \frac{z}{1 + \frac{z^2}{n}} \sqrt{\frac{\hat p(1 - \hat p)}{n} + \frac{z^2}{4n^2}},
  \qquad z = z_{1 - \alpha/2}.
\end{equation}
Unlike the Wald interval $\hat p \pm z\sqrt{\hat p(1-\hat p)/n}$, it stays inside $[0, 1]$ and keeps its coverage
near the nominal level when $p$ is close to the boundary, which is where many of the dataset's shares lie (null rates,
fraud rates, defect rates). The exact interval of \textcite{clopper1934} guarantees at least the nominal coverage at
the price of width; we use it only when a guarantee is required. For a count of $k$ events over exposure $E$, the
exact Poisson interval \autocite{garwood1936} is
\begin{equation}\label{eq:garwood}
  \left[\frac{\chi^2_{2k,\,\alpha/2}}{2E},\; \frac{\chi^2_{2k+2,\,1-\alpha/2}}{2E}\right],
\end{equation}
with the lower bound set to zero when $k = 0$. Every interval in this report has level $1 - \alpha =
\proto{intervals.level}$.

\section{Why independence fails, and the cluster bootstrap}\label{sec:em-clusters}

Suppose the outcomes of a cluster share a common component, so that any pair of observations of the same cluster
has correlation $\rho$, and all clusters have size $m$. Then
\begin{equation}\label{eq:deff}
  \Var\bigl(\hat\theta\bigr) = \frac{\sigma^2}{n}\bigl(1 + (m - 1)\rho\bigr),
\end{equation}
the design effect of \textcite{kish1965}. An interval that assumes $\rho = 0$ is too narrow by the factor
$\sqrt{1 + (m-1)\rho}$, which grows without bound with the cluster size. In this dataset clusters are large: a
customer has many transactions, an agent many contact hours. A small within-customer correlation is enough to make a
row-level interval useless. The fix does not need $\rho$ or a model of it: resample the independent units.

\begin{definition}[Cluster bootstrap]\label{def:cluster-bootstrap}
Draw $G$ clusters with replacement from $\{1, \dots, G\}$, keeping all rows of each drawn cluster, and recompute
$\hat\theta^{*} = \sum_{g^*} Y_{g^*} \big/ \sum_{g^*} D_{g^*}$. Repeat $B$ times. The percentile interval is
$[\hat\theta^{*}_{(\alpha/2)}, \hat\theta^{*}_{(1-\alpha/2)}]$, the empirical quantiles of the replicates.
\end{definition}

Under \cref{ass:clusters} the cluster totals are an independent sample, so the bootstrap of
\textcite{efron1979} applies to them directly; resampling clusters rather than rows is what keeps the within-cluster
dependence intact \autocite{davison1997,field2007}. Its large-sample justification is the same as the sandwich
variance of \textcite{liang1986}, which this bootstrap estimates without a parametric model, and it is reliable when
the number of clusters is large. The report uses it from $G = \proto{intervals.min_clusters}$ clusters upward. The test
suite checks the claim behind this section on simulated nested data with a known answer: the interval of
\cref{eq:wilson} covers the truth far less often than its nominal level, and the cluster bootstrap covers it close to
the nominal level.

\section{Few clusters}\label{sec:em-few}

With few clusters the sandwich variance is biased downwards and the bootstrap distribution is lumpy
\autocite{cameron2008,mackinnon2023}. Between $\proto{intervals.min_clusters_few}$ and
$\proto{intervals.min_clusters}$ clusters the report writes $\hat\theta$ as a weighted mean of cluster ratios,
$\hat\theta = \sum_g w_g r_g$ with $w_g = D_g / \sum_h D_h$ and $r_g = Y_g / D_g$, and corrects the variance as follows.

\paragraph{\CRtwo.} The bias-reduced variance of \textcite{bell2002} rescales each cluster's residual by its leverage.
For this intercept-only weighted model the leverage of cluster $g$ is $h_g = w_g$, and
\begin{equation}\label{eq:cr2}
  \widehat{\Var}_{\mathrm{\CRtwo}}\bigl(\hat\theta\bigr)
  = \sum_{g} w_g^2 \, \frac{(r_g - \hat\theta)^2}{1 - h_g},
\end{equation}
with the interval $\hat\theta \pm t_{G-1,\,1-\alpha/2}\,\widehat{\Var}_{\mathrm{\CRtwo}}^{1/2}$.

\paragraph{Wild cluster bootstrap.} Following \textcite{cameron2008}, each replicate multiplies every cluster's
residual by an independent Rademacher sign $v_g \in \{-1, 1\}$, sets $r_g^{*} = \hat\theta + v_g (r_g - \hat\theta)$,
and records the studentised statistic $t^{*} = (\hat\theta^{*} - \hat\theta) / \mathrm{se}_{\mathrm{\CRtwo}}^{*}$. The
interval inverts the quantiles of $t^{*}$ around the observed \CRtwo{} standard error. Below
$\proto{intervals.min_clusters_few}$ clusters, as with the markets of the dataset taken as clusters, no interval is
reported: the claim is descriptive and says why.

\section{Serial dependence}\label{sec:em-series}

Days are not exchangeable: volumes, contacts and errors on consecutive days are correlated, so resampling days
independently understates the variance of a mean over days for the same reason as \cref{eq:deff}. The report uses the
stationary bootstrap of \textcite{politis1994}: a replicate is a sequence of blocks of consecutive days, each block
starting at a uniform day and having a geometric length with mean $L$, the series being read as a circle. The mean
block length follows the automatic rule of \textcite{politis2004} as corrected by \textcite{patton2009}:
\begin{equation}\label{eq:pw}
  \hat b = \left(\frac{2\,\hat G^2}{\hat D}\right)^{1/3} n^{1/3},
  \qquad
  \hat G = \sum_{|k| \le M} \lambda(k/M)\,|k|\,\hat R(k),
  \qquad
  \hat D = 2\Bigl(\sum_{|k| \le M} \lambda(k/M)\,\hat R(k)\Bigr)^2,
\end{equation}
where $\hat R$ is the sample autocovariance, $\lambda$ the flat-top window that equals $1$ on $[0, 1/2]$ and falls
linearly to $0$ at $1$, and $M$ the bandwidth chosen from the first run of autocorrelations that stay within their
noise band. Applied to the daily differences that a gate averages, \cref{eq:pw} makes the block length grow with the
dependence the data show and shrink to single days when there is none.

\section{Many tests}\label{sec:em-multiplicity}

The notebooks ran many hypothesis tests. Read one at a time, a few small p-values among many are expected by chance,
so the report controls the false discovery rate, the expected share of false rejections among all rejections
\autocite{benjamini1995}. For ordered p-values $p_{(1)} \le \dots \le p_{(m)}$ of a family of $m$ tests, the
adjusted values are
\begin{equation}\label{eq:bh}
  q_{(i)} = \min_{j \ge i} \min\Bigl(1, \frac{m\,p_{(j)}}{j}\Bigr),
\end{equation}
and rejecting every hypothesis with $q \le \alpha$ controls the false discovery rate at $\alpha$ when the tests are
independent or positively dependent \autocite{benjamini2001}. The rate is only controlled over the family it is
computed on, so families are defined by where the tests come from, never by which tests the text happens to cite. The
protocol fixes them before any registered result exists: the notebooks' own declared families of the daily series
($\proto{families.daily-series-i.size}$ and $\proto{families.daily-series-ii.size}$ tests) and of the hourly series
($\proto{families.hourly-series-iii.size}$ tests), and the calendar effects of the country series on the main and the
backup copies ($\proto{families.country-calendar-main.size}$ and $\proto{families.country-calendar-backup.size}$
tests). The report's implementation of \cref{eq:bh} reproduces the adjusted values the notebooks stored for their
declared families, and agrees with independent implementations on every p-value column of the stored results.
A sentence that calls a result statistically detectable must cite its $q$.

\section{Specification search and the registered re-estimation}\label{sec:em-registered}

The analyses behind the Grain Atlas were exploratory: across many notebooks, choices of features, windows, benchmarks
and tests were made after looking at the data. Such choices can make chance findings look real
\autocite{simmons2011}. The standard remedy is to fix the analysis before seeing its result
\autocite{nosek2018}. This report does what is possible here and states what is not.

The dataset ends on \proto{dataset.end_date}, and every notebook saw data up to that day. There is no untouched
holdout. What can still be fixed in advance is the specification. Before any clustered estimate of a gate existed,
the protocol froze, in the repository's history, the gates to re-estimate (all $\proto{gates.ids.n}$ of them,
whatever their published verdict), each gate's out-of-time split as coded, its cluster unit, the verdict rule and the
rules that make a gate not assessable, the rule for comparing the main and backup copies, the test families, and the
digest of every data file. Any later change needs a dated amendment with a reason, and the build checks every frozen
field against the freezing commit.

A registered re-estimation therefore rules out adjusting the specification after seeing the clustered result, and an
interval that ignores the dependence between rows. It does not rule out that the original choices of features, windows
or benchmarks were shaped by the same data. It changes only how the uncertainty of a gate is measured, so a change of
verdict can be attributed to that and to nothing else.

\section{The readiness gate}\label{sec:em-gate}

A candidate model is trained out of time and compared with a transparent benchmark on the same test rows. Let
$\Delta$ be its gain (a higher AUC, a lower error or loss, expressed so that larger is better), $\Delta^{*}_1, \dots,
\Delta^{*}_B$ its bootstrap replicates, $[L, U]$ their $\alpha/2$ and $1 - \alpha/2$ quantiles with
$\alpha = \proto{gates.alpha}$, and $\mu$ the smallest gain worth deploying, set per model before training. The
minimum detectable effect at the test volume is
\begin{equation}\label{eq:mde}
  \mathrm{MDE} = c\,\widehat{\mathrm{sd}}\bigl(\Delta^{*}\bigr),
  \qquad c = \proto{gates.mde_multiplier} \approx z_{1-\alpha/2} + z_{\pi},\; \pi = \proto{gates.power}.
\end{equation}

\begin{definition}[Readiness verdict]\label{def:verdict}
The gate is \emph{green} when $L > 0$, $\Delta \ge \mu$ and, when folds are evaluated, every fold's gain has the
sign of $\Delta$; \emph{amber} when $L > 0$ but the gain is not material or the folds disagree; \emph{red} otherwise,
with the root cause \emph{insufficient volume} when $\mathrm{MDE} > \mu$ (even a material gain could not be told
from noise) and the scenario's structural cause when it is not.
\end{definition}

The registered re-estimation keeps \cref{def:verdict} unchanged and replaces only the replicates: each gate's
bootstrap resamples its independent unit (customers, agents, branches, campaigns, or stationary blocks of days) instead
of rows. It adds the following rules. A gate is \emph{not assessable} when it has fewer than $\proto{gates.min_clusters}$
clusters, because a cluster bootstrap with so few units cannot be trusted, and when more than a share
$\proto{gates.nan_share_max}$ of its replicates fail, for example because a resample of customers contains a single
class and its AUC is undefined. A re-estimated gate disagrees with the published scorecard when its verdict or its
root cause changes; \cref{ch:readiness-gates} reports every gate, on the main and the backup copies, under this rule.
